Embedding sharing cuts recommender memory, but can discard item-level popularity information and change which training objective performs best. Prior work connects embedding magnitudes to popularity and controls it through normalization or item-specific terms. Earlier compression methods use random or frequency-aware hashing; recent graph-based approaches group items by interaction structure. Neither line asks what sharing does to the popularity signal. In an empirical and analytical diagnosis, we identify a pooling bottleneck: shared embedding norms track bucket traffic (Spearman 0.93–0.94) more strongly than individual-item popularity (0.28–0.52). We trace a reversal to it: on Gowalla, DirectAU outperforms BPR with full tables (18.02\% versus 16.55\% Recall@20), but falls behind under graph-based sharing (7.55\% versus 9.37\%). We apply a correlation-ratio bound to quantify this representational constraint and use fixed-coefficient interventions to separate bucket calibration from item-specific information. An established, validation-tuned log-popularity correction, applied without retraining, raises compressed DirectAU to 13.30\% on Gowalla; item-specific information accounts for 71–76\% of the measured gain on two datasets. Corrected DirectAU exceeds equally corrected BPR by 0.83–2.13 percentage points across three datasets with three seeds each; all paired-user confidence intervals are positive. Stored in float32, the correction adds 3.9–7.3\% to compressed-table storage, while accuracy gains favor popular items: tail recall remains below full-table references and catalog coverage contracts.
